\note{Now that we have seen how AdvNet searches efficiently, let's look at how developers can use it to answer practical questions. By simply changing how we define the \textbf{reference} and the \textbf{target} inside the same scoring pipeline, the authors apply AdvNet to three distinct use cases:\par \par * \textbf{UC1 — Finding suboptimal performance and fairness issues}: We set the target to be the protocol under test, and compare it against either theoretical link capacity or another protocol implementation (or compare a reference flow running alone versus competing head-to-head with the target) to uncover severe underutilization, kernel bugs, or starvation.\par * \textbf{UC2 — Regression testing across versions}: We set the older version of a protocol as the reference and the newer version as the target, searching for environments where an update degrades throughput or fairness.\par * \textbf{UC3 — Validating design objectives}: To test the "do no harm" goal of multipath transport, we set a single-path configuration as the reference and a dual-path configuration as the target.\par \par Let's examine what AdvNet discovered in each of these three use cases.
\par\smallskip Source: AdvNet paper §1–3}
